\section{Related Work}
\label{s_rw}
\textbf{Model-Based Deformable Objects Manipulation}. To learn deformable object manipulation with unknown dynamics, previous works have modeled deformable objects using various state representations, such as particle systems~\cite{jaillet1998deformable,li2018learning,liu2023robotic}, points~\cite{sun2014heuristic,weng2022fabricflownet,yuba2017unfolding,jangir2020dynamic}, graphs~\cite{triantafyllou2016geometric,yuba2017unfolding,longhini2023edo,lippi2020latent}, and differentiable simulators~\cite{li2018learning,chen2024differentiable,heiden2021disect,schenck2018spnets,xian2023fluidlab}. They then learn or plan manipulation strategies using demonstrations~\cite{chen2024differentiable,balaguer2011combining}, annotations~\cite{chi2024iterative}, or simulation~\cite{li2018learning,lin2022learning}. However, because state estimation can fail under hand-eye coordination or object occlusions~\cite{chen2024differentiable}, generalizing this representation to handle unseen object dynamics, categories, and instances in the real world remains challenging. In contrast, \methodname{} is a model-free method that learns to infer a dynamics embedding from recent robot visual observations and rapidly adapts to unknown deformable object dynamics in the real world without requiring state estimation. This allows \methodname{} to perform real-robot mobile manipulation tasks that require hand-eye coordination and robustness against object occlusions, such as the \taskone~and~\tasktwo~mobile manipulation tasks in Sec.~\ref{sec:experiments}.

\textbf{Rapid Motor Adaptation.} Our goal is to achieve rapid adaptation to unknown dynamics of deformable objects, therefore we revise previous works that apply Rapid Motor Adaptation~\cite{kumar2021rma} to other robotic tasks, such as quadruped locomotion~\cite{fu2022deep,fu2021coupling,fu2021minimizing,zhuang2023robot}, bipedal locomotion~\cite{kumar2022adapting,lee2020learning}, in-hand manipulation~\cite{qi2023hand}, and simulated robotic manipulation~\cite{liang2024rapid}. These methods first identify the task-relevant dynamics properties, such as mass and friction, obtain and encode the values of these properties from the simulator into a dynamics embedding, and learn a rapid adaptation module to infer this embedding from robot observations and actions. However, effective deformable object manipulation requires not only rapid inference and adaptation to the object's mass and position, but also how the shape of the objects changes, an element missing from these methods. To address this challenge, we revise \methodname{} to infer a dynamics 
embedding of both the deformable object's mass and position and the shape changes of the object, from its ground-truth physics parameters and recent particle positions in simulation, respectively, and then learns to infer this dynamics embedding from robot visual observations and actions instead. This allows \methodname{} to learn a visuomotor policy in simulation, transfer it to a real-robot mobile manipulator, and manipulate real-world deformable objects using only onboard sensors. 

\textbf{Learning Deformable Object Manipulation from Real-World Data.}
To learn generalizable deformable object manipulation, prior methods have collected real-world data such as random trajectories~\cite{wang2019learning}, human demonstrations~\cite{seita2019deep,seita2021learning,joshi2017robotic}, human annotations~\cite{lui2013tangled}, and large-scale teleoperation data~\cite{lee2015learning} for both policy pre-training and fine-tuning~\cite{black2024pi0,kim2024openvla,team2024octo,zhao2024aloha}. In particular, the last family first bootstraps learning from internet-scale vision-language pre-training and crowd-sourced cross-embodiment robotics datasets~\cite{o2024open}, before fine-tuning on robot-specific datasets to improve target-embodiment performance. While such large-scale real-world data collection methods have demonstrated impressive task complexity, they currently do not zero-shot generalize to new robot embodiments such as the TIAGo mobile manipulator used in our experiments, and whether it can collect enough data to achieve such zero-shot generalization remains to be seen. As a less labor-intensive alternative to this family of algorithms, \methodname{} learns deformable object mobile manipulation purely in simulation instead and zero-shot transfers to a vision-based mobile manipulator while generalizing across unseen object dynamics, categories, and instances.

\begin{figure*}[t]
\centering
\includegraphics[width=\linewidth]{photos/objects.png}
\caption{All unseen real-world deformable objects used for the \taskone~(\textit{left}) and \tasktwo~(\textit{right}) mobile manipulation tasks, respectively. \textbf{1D categories (\textit{left}):} belt, hammer, marker, banana, HDMI adapter, GPU connector, jumping rope, ethernet cable, zip tie, rake, mouse, screwdriver, stick, rubber, velvet rope, VGA adapter, anchor rope, garden rope, ribbon, fan. \textbf{2D categories (\textit{right}):} polyester bag, pants, towel, plastic bag, face mask, paper bag, glove, shorts, cap, lid, sock, soft cuff, cardboard, envelope, pouch, wallet, eye-mask, t-shirt, sponge, wool hat. Learning from deformable objects across a wide range of instances in simulation enables \methodname{} to zero-shot transfer to real-world objects under unseen dynamics, instances, and categories.}
\label{fig:objects}
\vspace{-15px}
\end{figure*}
\textbf{Learning Deformable Object Manipulation from Simulation.}
To learn deformable object manipulation from simulation and transfer learned policies to the real world, previous works have adopted various strategies such as domain randomization~\cite{matas2018sim,hietala2021closing,ha2022flingbot,wu2019learning}, domain adaptation~\cite{chebotar2019closing}, real2sim~\cite{salhotra2022learning}, and dense object descriptors~\cite{sundaresan2020learning,ganapathi2021learning}. However, these methods do not include multiple visual observations and robot actions as input into their manipulation policy and therefore cannot infer and adapt to deformable object dynamics online. In particular, although Chebotar1 et al.~\cite{chebotar2019closing} proposes to adjust simulation parameters based on real-world experience, this method requires multiple real-world trajectories for adaptation training. It therefore cannot rapidly adapt to unknown dynamics in real-time. In comparison, \methodname{} leverages privileged simulator information to infer deformable object dynamics but then learns to infer them using robot visual observations and actions instead. This allows \methodname{} to rapidly adapt to a wide range of dynamics across real-world deformable objects, from soft objects (e.g., ribbons, cloths) to stiff items (e.g., cables, wallets).